% Part: normal-modal-logic
% Chapter: tableaux
% Section: rules-for-K

\documentclass[../../../include/open-logic-section]{subfiles}

\begin{document}

\olfileid[tr]{nml}{tab}{rul}

\olsection{\Ax{K} için Kurallar}

Olağan önerme bağlaçlarının kuralları, yalnızca önekler eklenmiş olmak
üzere, olağan önermesel işaretli !!{tableau}s kurallarıyla aynıdır. Her
durumda, $\sFmla{S}{!A}[\sigma]$ işaretli !!{formula}{}üne uygulanan kural,
yine $\sigma$ ile öneklenmiş yeni !!{formula}s üretir. Bu sezgisel olarak
açıktır: örneğin $!A \land !B$, $\sigma$'nın adlandırdığı dünyada doğruysa
$!A$ ile $!B$ de $\sigma$'da (başka bir dünyada değil) doğrudur. Önermesel
kuralları \olref{tab:prop-rules}'te bir araya getiriyoruz.

\begin{table}
  \[\def\arraystretch{3}\begin{array}{|c|c|}
    \hline
    \AxiomC{\sFmla{\True}{\lnot !A}[\sigma]}
    \RightLabel{\TRule{\True}{\lnot}}
    \UnaryInfC{\sFmla{\False}{!A}[\sigma]}
    \DisplayProof
    &
    \AxiomC{\sFmla{\False}{\lnot !A}[\sigma]}
    \RightLabel{\TRule{\False}{\lnot}}
    \UnaryInfC{\sFmla{\True}{!A}[\sigma]}
    \DisplayProof
    \\[1ex]
    \hline
    \AxiomC{\sFmla{\True}{!A \land !B}[\sigma]}
    \RightLabel{\TRule{\True}{\land}}
    \UnaryInfC{\sFmla{\True}{!A}[\sigma]}
    \noLine
    \UnaryInfC{\sFmla{\True}{!B}[\sigma]}
    \DisplayProof
    &
    \AxiomC{\sFmla{\False}{!A \land !B}[\sigma]}
    \RightLabel{\TRule{\False}{\land}}
    \UnaryInfC{$\sFmla{\False}{!A}[\sigma] \quad \mid \quad
      \sFmla{\False}{!B}[\sigma]$}
    \DisplayProof
    \\[2ex]
    \hline
    \AxiomC{\sFmla{\True}{!A \lor !B}[\sigma]}
    \RightLabel{\TRule{\True}{\lor}}
    \UnaryInfC{$\sFmla{\True}{!A}[\sigma] \quad \mid \quad
      \sFmla{\True}{!B}[\sigma]$}
    \DisplayProof
    &
    \AxiomC{\sFmla{\False}{!A \lor !B}[\sigma]}
    \RightLabel{\TRule{\False}{\lor}}
    \UnaryInfC{\sFmla{\False}{!A}[\sigma]}
    \noLine
    \UnaryInfC{\sFmla{\False}{!B}[\sigma]}
    \DisplayProof
    \\[2ex]
    \hline
    \AxiomC{\sFmla{\True}{!A \lif !B}[\sigma]}
    \RightLabel{\TRule{\True}{\lif}}
    \UnaryInfC{$\sFmla{\False}{!A}[\sigma] \quad \mid
      \quad \sFmla{\True}{!B}[\sigma]$}
    \DisplayProof
    &
    \AxiomC{\sFmla{\False}{!A \lif !B}[\sigma]}
    \RightLabel{\TRule{\False}{\lif}}
    \UnaryInfC{\sFmla{\True}{!A}[\sigma]}
    \noLine
    \UnaryInfC{\sFmla{\False}{!B}[\sigma]}
    \DisplayProof
    \\[2ex]
    \hline
  \end{array}\]
  \caption{Önerme bağlaçları için önekli !!{tableau} kuralları}
  \ollabel{tab:prop-rules}
\end{table}

Kapanma koşulu olağan !!{tableau}s için olanla aynıdır; ancak yalnızca
!!{formula}s{}in değil, öneklerin de eşleşmesini isteriz. Dolayısıyla bir dal,
bir $\sigma$ öneki ve bir !!{formula} $!A$ için hem
\[
\sFmla{\True}{!A}[\sigma] \quad\text{hem}\quad \sFmla{\False}{!A}[\sigma]
\]
içeriyorsa kapalıdır.

Varsayımları yerleştirme kuralları da olağan !!{tableau}s için olanla aynıdır;
tek fark, varsayımlar için her zaman $1$ önekini kullanmamızdır. (Bütün
varsayımlarda aynı öneki kullandığımız sürece hangisini seçtiğimiz önemli
değildir.) Örneğin
\[
!B_1, \dots, !B_n \Proves !A
\]
ancak ve ancak şu varsayımlar için kapalı bir tablo varsa geçerlidir:
\[
\sFmla{\True}{!B_1}[1], \dots, \sFmla{\True}{!B_n}[1],
\sFmla{\False}{!A}[1].
\]

\iftag{prvBox}{\iftag{prvDiamond}{$\Box$ ve
    }{$\Box$}}{}\iftag{prvDiamond}{$\Diamond$}{} biçimindeki modal işleçler için,
$\sigma$ önekli !!a{formula}{}e uygulanan kuralın sonucunun öneki $\sigma.n$
olur. Ancak hangi $n$ değerine izin verildiği, işaretin $\True$ ya da
$\False$ olmasına bağlıdır.

\iftag{prvBox}{$\TRule{\True}{\Box}$ kuralı,
  $\sFmla{\True}{\Box !A}[\sigma]$ içeren bir dalı
  $\sFmla{\True}{!A}[\sigma.n]$ ile genişletir.\iftag{prvDiamond}{ Benzer
    biçimde, }{}}{}\iftag{prvDiamond}{$\TRule{\False}{\Diamond}$ kuralı,
  $\sFmla{\False}{\Diamond !A}[\sigma]$ içeren bir dalı
  $\sFmla{\False}{!A}[\sigma.n]$ ile genişletir.}{}
\iftag{notprvBox,notprvDiamond}{Bu kural}{Bu kurallar}, yalnızca uygulandığı
dalda \emph{zaten} bulunan bir $\sigma.n$ öneki için uygulanabilir. Böyle bir
öneke dalda ``kullanılmış'' diyelim.

\iftag{prvBox}{$\TRule{\False}{\Box}$ kuralı,
  $\sFmla{\False}{\Box !A}[\sigma]$ içeren bir dalı
  $\sFmla{\False}{!A}[\sigma.n]$ ile genişletir.\iftag{prvDiamond}{ Benzer
    biçimde, }{}}{}\iftag{prvDiamond}{$\TRule{\True}{\Diamond}$ kuralı,
  $\sFmla{\True}{\Diamond !A}[\sigma]$ içeren bir dalı
  $\sFmla{\True}{!A}[\sigma.n]$ ile genişletir.}{}
Ne var ki \iftag{notprvBox,notprvDiamond}{bu kural}{bu kurallar}, yalnızca
uygulandığı dalda \emph{henüz} bulunmayan bir $\sigma.n$ öneki için
uygulanabilir. Böyle öneklere dala göre ``yeni'' deriz.

Kurallar \olref{tab:rules-K}'de verilmiştir.

\begin{table}
  \begin{center}
    \def\arraystretch{3}\def\fCenter{}
    \begin{tabular}{|c|c|}
    \hline
    \iftag{prvBox}{
      \AxiomC{\sFmla{\True}{\Box !A}[\sigma]}
      \RightLabel{\TRule{\True}{\Box}}
      \UnaryInfC{\sFmla{\True}{!A}[\sigma.n]}
      \DisplayProof
      &
      \AxiomC{\sFmla{\False}{\Box !A}[\sigma]}
      \RightLabel{\TRule{\False}{\Box}}
      \UnaryInfC{\sFmla{\False}{!A}[\sigma.n]}
      \DisplayProof\\
      $\sigma.n$ kullanılmıştır & $\sigma.n$ yenidir
      \\[1ex]
      \hline}{}
    \iftag{prvDiamond}{
      \AxiomC{\sFmla{\True}{\Diamond !A}[\sigma]}
      \RightLabel{\TRule{\True}{\Diamond}}
      \UnaryInfC{\sFmla{\True}{!A}[\sigma.n]}
      \DisplayProof
      &
      \AxiomC{\sFmla{\False}{\Diamond !A}[\sigma]}
      \RightLabel{\TRule{\False}{\Diamond}}
      \UnaryInfC{\sFmla{\False}{!A}[\sigma.n]}
      \DisplayProof\\
      $\sigma.n$ yenidir & $\sigma.n$ kullanılmıştır
      \\[1ex]
      \hline}{}
    \end{tabular}
  \end{center}
  \caption{\Ax{K} için modal kurallar.}
  \ollabel{tab:rules-K}
\end{table}

\iftag{prvBox}{\TRule{\True}{\Box}}{\TRule{\False}{\Diamond}} için
önekin kullanılmış olmasını şart koşmak gerekir; aksi hâlde aşağıdakini kapalı
bir !!{tableau} sayardık:
\iftag{notprvBox,notprvDiamond}{%
  \iftag{prvBox}{%
  \begin{oltableau}
    [\pFmla{\True}{\Box \formula{A}}{1}, just = \TAss
      [\pFmla{\False}{\lnot\Box\lnot \formula{A}}{1}, just = \TAss
        [\pFmla{\True}{\formula{A}}{1.1}, just = {\TRule{\True}{\Box}[1]}
          [\pFmla{\True}{\Box\lnot\formula{A}}{1},
            just ={\TRule{\False}{\lnot}[2]}
            [\pFmla{\True}{\lnot\formula{A}}{1.1},
              just = {\TRule{\True}{\Box}[4]}
              [\pFmla{\False}{\formula{A}}{1.1},
                  just ={\TRule{\True}{\lnot}[5]}, close]
            ]
          ]
        ]
      ]
    ]
  \end{oltableau}
  }{
  \begin{oltableau}
    [\pFmla{\True}{\lnot\Diamond\lnot \formula{A}}{1}, just = \TAss
      [\pFmla{\False}{\Diamond\formula{A}}{1}, just = \TAss
        [\pFmla{\False}{\formula{A}}{1.1}, just = {\TRule{\False}{\Diamond}[2]}
          [\pFmla{\False}{\Diamond\lnot\formula{A}}{1},
            just ={\TRule{\True}{\lnot}[1]}
            [\pFmla{\False}{\lnot\formula{A}}{1.1},
              just = {\TRule{\False}{\Diamond}[4]}
              [\pFmla{\True}{\formula{A}}{1.1},
                  just ={\TRule{\True}{\lnot}[5]}, close]
            ]
          ]
        ]
      ]
    ]
  \end{oltableau}
}}{
  \begin{oltableau}
    [\pFmla{\True}{\Box \formula{A}}{1}, just = \TAss
      [\pFmla{\False}{\Diamond \formula{A}}{1}, just = \TAss
        [\pFmla{\True}{\formula{A}}{1.1}, just = {\TRule{\True}{\Box}[1]}
          [\pFmla{\False}{\formula{A}}{1.1},
            just = {\TRule{\False}{\Diamond}[2]}, close]
        ]
      ]
    ]
\end{oltableau}}

Fakat $\Box \formula{A} \Entails/ \Diamond \formula{A}$; dolayısıyla ispat
sistemimiz sağlam olmazdı. Benzer biçimde,
$\Diamond \formula{A} \Entails/ \Box \formula{A}$; ancak
\iftag{prvBox}{\TRule{\False}{\Box}}{\TRule{\True}{\Diamond}} için
önekin yeni olması kısıtı bulunmasaydı aşağıdaki tablo kapalı olurdu:
\iftag{notprvBox,notprvDiamond}{%
  \iftag{prvBox}{%
  \begin{oltableau}
    [\pFmla{\True}{\lnot\Box\lnot \formula{A}}{1}, just = \TAss
      [\pFmla{\False}{\Box \formula{A}}{1}, just = \TAss
        [\pFmla{\False}{\formula{A}}{1.1}, just = {\TRule{\True}{\Box}[2]}
          [\pFmla{\False}{\Box\lnot\formula{A}}{1},
            just ={\TRule{\True}{\lnot}[1]}
            [\pFmla{\False}{\lnot\formula{A}}{1.1},
              just = {\TRule{\False}{\Box}[4]}
              [\pFmla{\True}{\formula{A}}{1.1},
                  just ={\TRule{\False}{\lnot}[5]}, close]
            ]
          ]
        ]
      ]
    ]
  \end{oltableau}
  }{
  \begin{oltableau}
    [\pFmla{\True}{\Diamond \formula{A}}{1}, just = \TAss
      [\pFmla{\False}{\lnot\Diamond\lnot\formula{A}}{1}, just = \TAss
        [\pFmla{\True}{\formula{A}}{1.1}, just = {\TRule{\True}{\Diamond}[1]}
          [\pFmla{\True}{\Diamond\lnot\formula{A}}{1},
            just ={\TRule{\False}{\lnot}[2]}
            [\pFmla{\True}{\lnot\formula{A}}{1.1},
              just = {\TRule{\True}{\Diamond}[4]}
              [\pFmla{\False}{\formula{A}}{1.1},
                  just ={\TRule{\True}{\lnot}[5]}, close]
            ]
          ]
        ]
      ]
    ]
  \end{oltableau}
}}{
  \begin{oltableau}
    [\pFmla{\True}{\Diamond \formula{A}}{1}, just = \TAss,name=one
      [\pFmla{\False}{\Box \formula{A}}{1}, just = \TAss, name=two
        [\pFmla{\True}{\formula{A}}{1.1}, just = {\TRule{\True}{\Diamond}[1]}
          [\pFmla{\False}{\formula{A}}{1.1},
            just = {\TRule{\False}{\Box}[2]}, close]
        ]
      ]
    ]
  \end{oltableau}}

\end{document}
